\documentclass[12pt]{article}
\usepackage[T1]{fontenc}
\usepackage[utf8]{inputenc}
\usepackage{lmodern}
\usepackage{textcomp}
\usepackage{enumitem}
\usepackage{geometry}
\geometry{margin=1in}
\begin{document}
\begin{center}
Answer Key - Psychology - Graduate Level Assessment 5 \
\small (Answers are ordered exactly as in the question paper.)
\end{center}

\section*{Multiple Choice}
\begin{enumerate}[label=\arabic*.]
\item \textbf{Answer:} A
\\ \textit{Options:} A) Wechsler Intelligence Scale for Children-IV; B) Kaufman Assessment Battery for Children-II; C) Cognitive Assessment System-II; D) Woodcock-Johnson IV Tests of Cognitive Abilities; E) Differential Ability Scales-II
\\ \textit{Note:} The Wechsler Intelligence Scale for Children-IV is designed to assess cognitive abilities in children and adolescents, including those who are nonverbal, making it the most appropriate choice for evaluating general cognitive abilities in this population.
\item \textbf{Answer:} A
\\ \textit{Options:} A) Their neurons continue to fire at a rapid pace, keeping them in a state of high alert.; B) Their neurons remain in a state of graded potentiality even after they have fired.; C) Their heart rate remains elevated due to the stress of the crisis.; D) Their muscles remain tense due to the physical exertion of the crisis.; E) Their pituitary gland continues to release stress hormones.
\\ \textit{Note:} The correct answer highlights the neurological response to crises, where heightened neuronal activity maintains a state of arousal and vigilance, essential for survival and coping with stress.
\item \textbf{Answer:} A
\\ \textit{Options:} A) Broad stratum abilities; B) Crystallized intelligence; C) Abstract reasoning; D) Achievement or aptitude; E) Emotional intelligence
\\ \textit{Note:} Broad stratum abilities encompass a range of cognitive skills that are less influenced by cultural factors and are primarily assessed through nonverbal measures, highlighting their noncultural bias in evaluating intelligence.
\item \textbf{Answer:} A
\\ \textit{Options:} A) overattributing to person-situation interactions; B) overattributing to personal characteristics; C) overattributing to situations; D) overattributing to unpredictable events; E) underattributing to predictable events
\\ \textit{Note:} Attribution theorists identify the tendency to overattribute behavior to personal traits while underestimating situational influences as a key error in causal attributions, highlighting the importance of context in understanding others' actions.
\item \textbf{Answer:} D
\\ \textit{Options:} A) negative effects on productivity but little or no impact on attitudes.; B) positive effects on productivity but negative effects on attitudes.; C) long-lasting negative effects on attitudes and productivity.; D) positive effects on attitudes but little or no impact on productivity.; E) positive effects on productivity but little or no impact on attitudes.
\\ \textit{Note:} Research indicates that a compressed workweek can enhance employee morale and job satisfaction, leading to more positive attitudes, while evidence suggests that productivity levels may remain unchanged or show minimal variation due to the shorter work hours.
\end{enumerate}

\section*{True / False}
\begin{enumerate}[label=\arabic*.]
\setcounter{enumi}{5}
\item \textbf{Answer:} True
\\ \textit{Options:} A) True; B) False
\\ \textit{Note:} Vicarious liability is a significant concern for psychologists in supervisory roles because they may be held legally responsible for the actions and professional conduct of their trainees or staff, particularly if those actions result in harm to clients or violate ethical standards.
\item \textbf{Answer:} False
\\ \textit{Options:} A) True; B) False
\\ \textit{Note:} High shrinkage in multiple correlation coefficients indicates that the model's predictive power is diminished due to overfitting or the model not generalizing well to new data, rather than suggesting high reliability of the criterion variable.
\item \textbf{Answer:} True
\\ \textit{Options:} A) True; B) False
\\ \textit{Note:} Adler's individual psychology emphasizes that children's misbehavior is often motivated by underlying psychological goals, such as the desire for attention, the need for revenge, or the assertion of power, reflecting their attempts to navigate social relationships and fulfill their emotional needs.
\item \textbf{Answer:} False
\\ \textit{Options:} A) True; B) False
\\ \textit{Note:} While cognition, emotion, behavior, and environment are important factors in learning, effective learning outcomes also depend on additional elements such as motivation, prior knowledge, and instructional methods, making the statement incomplete and therefore false.
\item \textbf{Answer:} True
\\ \textit{Options:} A) True; B) False
\\ \textit{Note:} Research in psychology indicates that heightened stress levels can impair cognitive functioning and attention, increasing the likelihood of accidents in the workplace.
\end{enumerate}

\section*{Long Form Response}
\begin{enumerate}[label=\arabic*.]
\setcounter{enumi}{10}
\item \textbf{Answer:} Individuals low in self-monitoring often rely on their life scripts-internalized narratives shaped by personal experiences and societal expectations-to navigate social situations. These life scripts provide a framework for understanding appropriate behaviors and responses in various contexts, allowing low self-monitors to engage with others based on their established beliefs and values rather than adapting to social cues. This reliance on personal narratives can lead to a more authentic expression of self, but may also result in social missteps when the individual's script does not align with the expectations of the social environment. Consequently, while such reliance fosters a sense of identity and consistency, it can limit the flexibility needed to effectively navigate complex social interactions, impacting relational dynamics and overall social competence.
\\ \textit{Note:} correct because it highlights how individuals low in self-monitoring depend on their life scripts, which are shaped by personal experiences and societal norms, to navigate social interactions authentically, rather than adjusting their behavior based on external social cues.
\item \textbf{Answer:} Asking a candidate's age during personnel selection raises significant ethical considerations and psychological implications, particularly when age is considered a bona fide occupational requirement. While such a requirement may be justified in specific contexts-such as roles involving age-related expertise or legal regulations-the practice can perpetuate age discrimination and reinforce stereotypes about older or younger workers. Ethically, employers must balance the legitimate needs of the organization with the potential for bias against candidates based solely on age. Psychologically, this inquiry can lead to feelings of inadequacy and anxiety among candidates, affecting their self-esteem and willingness to engage in the application process. Therefore, while age may sometimes be a legitimate factor, it is crucial to implement selection criteria that uphold fairness and inclusivity, avoiding unnecessary emphasis on age that may detract from a candidate's qualifications and capabilities.
\\ \textit{Note:} correct because it addresses the ethical dilemmas of potential age discrimination and the reinforcement of stereotypes while recognizing that age can be a legitimate requirement in certain job contexts, thus highlighting the need for a careful balance in personnel selection practices.
\end{enumerate}

\vfill
\noindent\textit{End of Answer Key}
\end{document}